\documentclass[sigconf,nonacm]{acmart}
%% ACM SIGKDD Explorations target. Swap to the Explorations house style for camera-ready.
\usepackage{booktabs}
\usepackage{forest}
\usepackage{graphicx}

\begin{document}

\title{Fairness and Equity in LLM-Based Agents: A Taxonomic Survey}

\author{Rohith}
\authornote{Corresponding author.}
\affiliation{\institution{}\country{}}
\email{rohithreddybc98@gmail.com}

\author{Wenbin Zhang}
\affiliation{\institution{Florida International University}\country{USA}}

\begin{abstract}
Large language model (LLM) systems have crossed a threshold: they no longer merely \emph{answer}, they \emph{act}. Contemporary agents select tools and APIs, accumulate memory across turns, retrieve external evidence, delegate subtasks across roles, and plan over long horizons \cite{wang2024survey, xi2023rise, yao2023react}. These capabilities are now being wired into hiring, lending, healthcare, and public-benefits pipelines, where a biased \emph{decision} — not merely a biased sentence — has material consequences. Yet fairness research remains overwhelmingly anchored to single-turn question answering, scoring the words a model emits rather than the actions an agent takes. We present the first dedicated survey of fairness and equity in LLM-based agents, organized \textbf{by agent component} — the locus where bias enters — and cross-cut by an evaluation axis (group, individual, counterfactual; CFR/MASD; action-level disparity) and a mitigation axis (prompting, alignment, guardrails, multi-agent, process-level). We surface a structural \emph{measurement gap}: today's benchmarks grade answers, not actions. We corroborate it with a small original cross-framework bias audit exposing component-level disparity, and lay out a research agenda toward benchmarks for \emph{acting} agents.
\end{abstract}

\keywords{fairness, equity, large language models, LLM agents, bias, counterfactual
fairness, multi-agent systems, trustworthy AI}

\maketitle

\section{1. Introduction}

\subsection{1.1 From Answering to Acting}

The dominant mode of LLM deployment is shifting from \emph{generation} to \emph{agency}. An agent wraps a language model in a loop that perceives state, invokes tools, reads and writes memory, and commits to actions in an external environment \cite{wang2024survey, xi2023rise}. The reasoning-and-acting paradigm formalized by ReAct interleaves chain-of-thought with tool calls \cite{yao2023react}; cognitive-architecture accounts generalize this to perception, memory, planning, and action modules \cite{sumers2024cognitive}; and multi-agent frameworks compose several such loops into delegated, role-structured pipelines \cite{wu2023autogen, park2023generative}. The practical consequence is that an LLM increasingly does not return a paragraph for a human to act on — it \emph{selects the applicant to interview, retrieves the precedent that frames the triage, escalates or denies the claim.}

This shift raises the fairness stakes along three axes that single-turn QA does not exhibit. First, \textbf{autonomy}: an agent's output is a decision or an action with downstream effect, not a suggestion a human filters. Second, \textbf{compounding}: bias introduced at one component (a skewed retrieval, a stereotyped role assignment) propagates through subsequent steps and can amplify over a trajectory rather than washing out \cite{ma2026implicit, nguyen2025social}. Third, \textbf{opacity}: harm is distributed across tool calls, memory writes, and inter-agent messages, so a disparity visible in the final action may have no single attributable cause. These properties matter precisely in the \emph{consequential, autonomy-bearing} domains now adopting agents — hiring \cite{an2025measuring, veldanda2023emily, nghiem2024you}, lending and credit \cite{bowen2024measuring, feng2023empowering, azime2025accept}, clinical decision support \cite{omar2025sociodemographic, zack2024coding, poulain2024bias}, and public-benefits adjudication. Audits in these settings already document substantial disparities: resume evaluation exhibits intersectional gender-and-race penalties \cite{an2025measuring}; a mortgage-underwriting experiment finds racial disparities in LLM approval recommendations \cite{bowen2024measuring}; and sociodemographic cues shift medical management plans even when clinically irrelevant \cite{omar2025sociodemographic}. A discriminatory autonomous agent in these pipelines is a matter of national importance: it can systematically and silently allocate opportunity, credit, and care.

\subsection{1.2 Thesis: Fairness Is a Per-Component Property}

The organizing claim of this survey is that \textbf{fairness in an LLM agent cannot be captured by a single model-level metric; it must be analyzed per agent component}, because bias enters at distinct, separately-measurable, and separately-mitigable points in the pipeline. A model that is near-parity on a QA bias benchmark can still select tools unevenly across demographic framings \cite{blankenstein2025biasbusters}, inherit unfairness from a retrieval step that the base model alone never exhibited \cite{wu2024rag, hu2024free}, acquire bias from an assigned persona or role \cite{gupta2023bias, cao2026from}, or drift toward sycophantic, identity-conforming behavior over a long dialogue \cite{liu2025truth, ma2026implicit}. Crucially, \emph{agent decisions reveal implicit biases that direct questioning does not surface} \cite{li2025actions}: probing the answer misses the disparity in the act. Organizing the literature by component — rather than by fairness metric — is therefore not a presentational choice but the survey's central methodological contribution.

\subsection{1.3 Contributions}

1. \textbf{An agent-component taxonomy of where unfairness enters.} We map the literature onto the agent pipeline along six entry points — tool/API selection, memory and retrieval, multi-agent delegation and role assignment, planning and decomposition, user modeling and personalization, and long-horizon drift — plus a cross-cutting analysis of how component harms \emph{compound} (the agentic multiplier). For each, we give a mechanism of harm and representative evidence (Figure 1, Table 1).

2. \textbf{An evaluation framework for agentic fairness.} We map group, individual, and counterfactual fairness onto agent components, generalize the counterfactual flip rate (CFR) and mean absolute score difference (MASD) — formulated for a single customer-service agent by \cite{mayilvaghanan2026counterfactual} — into component- and trajectory-level instruments, and articulate \textbf{action-level disparity} metrics (decision, tool-invocation, and escalation disparity) that current benchmarks do not capture.

3. \textbf{A mitigation map} relating prompting, fine-tuning and alignment, guardrails, multi-agent debiasing, and process-level architectural interventions to the components they act on and the pipeline stage at which they intervene (Table 3).

4. \textbf{An original cross-framework bias audit} that exposes the measurement gap empirically: running multiple agent configurations on a consequential-decision task with counterfactual demographic swaps, we surface component-level action disparity that QA-level scoring of the same underlying model fails to register.

5. \textbf{An open-problems agenda}, culminating in the case for a fairness benchmark that scores \emph{acting} agents — covering long-horizon and compounding fairness, multi-agent dynamics, intersectional deployment harms, and agent-level (rather than token-level) mitigation.

\subsection{1.4 Positioning}

This survey is the \textbf{agentic successor to the parent QA-fairness lineage}. Chu et al.'s \emph{Fairness in Large Language Models: A Taxonomic Survey} \cite{chu2024fairness}, also in this venue, taxonomizes bias in LLMs that \emph{answer} — its axes are intrinsic/extrinsic bias, definitions, and debiasing of generation; it does not treat the agentic surface (tools, memory, delegation, planning, horizons) where, we argue, distinct fairness phenomena arise. We inherit its definitional grounding \cite{dwork2012fairness, hardt2016equality, kusner2017counterfactual, gallegos2024bias} and extend it to systems that act.

We differentiate, explicitly and honestly, from three near-neighbors. \emph{Evaluation and Benchmarking of LLM Agents: A Survey} \cite{mohammadi2025evaluation} is a broad agent-evaluation survey in which fairness is \textbf{one sub-topic among many}, not the organizing axis; we make it central and component-structured. \emph{Agentic AI in Healthcare & Medicine: A Seven-Dimensional Taxonomy} \cite{vatsal2026agentic} treats fairness as \textbf{one of seven dimensions} in a conceptual, single-domain framework; ours is fairness-dedicated, cross-domain, and tied to executable metrics. \emph{Counterfactual Fairness Evaluation of LLM-Based Contact Center Agent Quality Assurance} \cite{mayilvaghanan2026counterfactual} applies CFR/MASD in a \textbf{single domain} (customer service) and a single agent role; we generalize its method across every agent component and pair it with action-level disparity. To our knowledge, no \textbf{dedicated, agent-component-organized fairness-of-agents survey} exists as of mid-2026. We are careful here: the field is not unaware of fairness in agents — multi-agent bias \cite{borah2024implicit, madigan2025emergent}, trustworthy-agent threat surveys \cite{yu2025survey}, and component-specific audits \cite{blankenstein2025biasbusters, wu2024rag} all exist — but none assembles these into a component-organized taxonomy of fairness across the agent pipeline. We claim the \emph{first dedicated taxonomy}, not the \emph{first mention}. The remainder of the paper proceeds from background and scope (§2) to the core taxonomy (§3), evaluation (§4), mitigation (§5), our empirical audit (§6), a coverage cross-analysis (§7), and the research agenda (§8).



\section{2. Background & Scope}

\subsection{2.1 Fairness Definitions}

Formal fairness research has converged on three partially complementary families of criteria, each encoding a distinct moral intuition about equitable treatment.

\textbf{Group fairness} requires that outcomes be statistically equivalent across protected demographic groups. Dwork et al.'s foundational framework \cite{dwork2012fairness} formalizes \emph{fairness through awareness}, stipulating that similar individuals—judged by a task-relevant metric—receive similar outcomes, while simultaneously demanding that aggregate outcomes not be systematically skewed by group membership. Hardt et al. extend this to the supervised learning setting through \emph{equality of opportunity} and \emph{equalized odds}, requiring that true-positive and false-positive rates be equalized across groups conditioned on the true outcome label \cite{hardt2016equality}. These criteria are operationally tractable: they reduce fairness to a statistical comparison over a finite set of labeled groups and a classifier's output distribution.

\textbf{Individual fairness}, as articulated by Dwork et al. \cite{dwork2012fairness}, demands that similar individuals receive similar outcomes regardless of group identity, instantiating the classical liberal intuition that like cases be treated alike. In practice this requires specifying a task-appropriate similarity metric—a notoriously difficult design choice—which has limited the criterion's empirical adoption relative to group-level metrics.

\textbf{Counterfactual fairness} shifts the framing from outcomes to causal mechanisms \cite{kusner2017counterfactual}. Under this criterion, a decision is fair toward an individual if it would remain the same in a counterfactual world in which that individual's protected attributes had taken different values, holding all causally upstream factors constant. The criterion imports causal graph machinery into fairness reasoning: a prediction is deemed biased only if protected attributes appear in any causal path influencing that prediction. This causal perspective is particularly important for agentic settings, where multi-step decision chains create long causal paths along which demographic signals can propagate and amplify.

A comprehensive taxonomy of bias and fairness notions in classical machine learning is given by Mehrabi et al. \cite{mehrabi2021survey}, who catalog over twenty distinct bias sources spanning pre-processing, in-processing, and post-processing stages of the ML pipeline. That taxonomy serves as our baseline vocabulary; the present survey extends it upstream—into the agent components that condition what data reaches the model, how reasoning is structured, and how outputs are translated into consequential actions.

\subsection{2.2 LLM Bias: A Brief Inventory and the Agentic Gap}

The past four years have produced a rich literature measuring bias in large language models. Gallegos et al. provide an authoritative survey of bias and fairness across the LLM lifecycle, covering stereotyping, toxicity, representational harm, and distributional disparities in generated text \cite{gallegos2024bias}. Chu et al. offer a taxonomic treatment organized by fairness dimension and LLM task type, serving as the direct parent lineage for the present work \cite{chu2024fairness}.

Empirical measurement has been operationalized through a family of benchmark datasets that probe for demographic bias in model outputs. BBQ \cite{parrish2022bbq} uses ambiguous question-answering scenarios to surface stereotypical associations across nine social dimensions including age, race, religion, and disability status. BOLD \cite{dhamala2021bold} measures sentiment and regard disparities in open-ended generation by conditioning prompts on demographic group membership. StereoSet \cite{nadeem2021stereoset} evaluates the tradeoff between language modeling ability and stereotypical bias using intrasentence and intersentence contextual tasks. CrowS-Pairs \cite{nangia2020crows} presents masked-language-model pairs designed to surface biases favoring historically advantaged groups. The HolisticBias descriptor dataset \cite{smith2022sorry} broadens coverage to over 600 demographic descriptors across thirteen axes, enabling systematic measurement of sentiment and regard patterns in model responses.

This body of work has established that LLMs encode and reproduce social stereotypes, exhibit differential performance across demographic groups, and generate text with systematically unequal sentiment toward protected categories. However, all of these benchmarks evaluate \textbf{single-turn language generation}: a prompt is supplied, a response is generated, and a fairness metric is computed over the token distribution. This measurement paradigm is categorically insufficient for agentic systems, for at least three structural reasons.

First, agents act sequentially: a demographic signal present in a user's profile may influence an early tool selection, which constrains subsequent information retrieval, which in turn shapes a planning step that produces the final consequential decision. The bias does not appear in any single response but accumulates across the trajectory. Second, agents operate with persistent state: memory modules, retrieved context, and conversation history create channels through which biased inferences can carry forward across sessions and users—a mechanism absent from single-turn evaluation. Third, agents interact with external APIs and tools that may themselves encode biases in their outputs, creating a new source of harm not attributable to the LLM's parameters but to the LLM's tool-selection policy. BBQ can detect that a model answers a stereotyped question incorrectly; it cannot detect that a hiring agent systematically routes applicants with certain name patterns to lower-quality resume-screening tools. This measurement gap—the central observation of this survey—motivates the component-organized taxonomy developed in Section 3.

\subsection{2.3 What Is an LLM-Based Agent? Component Anatomy}

An LLM-based agent is a system in which a large language model serves as the core reasoning engine within an architecture that supports autonomous, multi-step interaction with external environments. We adopt the standard definition from the agent survey literature \cite{wang2024survey,xi2023rise}: an agent perceives inputs, reasons over its current context, selects and executes actions (including calling external tools), manages persistent memory, and iterates until a goal is achieved or a stopping criterion is met. The cognitive architectures survey by Sumers et al. provides a useful analytical framework for decomposing these systems into their functional constituents \cite{sumers2024cognitive}.

We identify six components whose design choices are consequential for fairness:

\textbf{1. Tool and API selection.} Agents use tool-calling interfaces to interact with external systems—databases, search engines, APIs, calculators, code executors. Schick et al.'s Toolformer demonstrates that LLMs can learn to invoke APIs autonomously by self-supervising on which calls improve language modeling performance \cite{schick2023toolformer}. The ReAct framework by Yao et al. interleaves reasoning traces with action steps, showing that the choice of which tool to call at each step is itself a learned behavior conditioned on contextual signals—including potentially demographic signals in the user query or retrieved context \cite{yao2023react}. Tool-selection policy is therefore a locus of potential discriminatory routing.

\textbf{2. Memory and retrieval.} Agents maintain both in-context history and external memory stores indexed for retrieval. Lewis et al. introduced retrieval-augmented generation (RAG) \cite{lewis2020retrieval}, in which relevant passages are retrieved from a document corpus and prepended to the model's context before generation. RAG architectures create a bias pathway distinct from the LLM's parameters: the retrieval function's ranking of documents can encode demographic disparities, and the model then generates conditioned on a potentially biased information set. Shinn et al.'s Reflexion \cite{shinn2023reflexion} demonstrates agents that accumulate verbal reflections as episodic memory; this same accumulation mechanism can encode and reinforce early biased inferences across subsequent reasoning steps.

\textbf{3. Multi-agent orchestration and delegation.} Modern deployments increasingly distribute tasks across networks of specialized agents. Wu et al.'s AutoGen framework \cite{wu2023autogen} enables hierarchical and lateral delegation between agents with distinct roles; Park et al.'s generative agent simulations \cite{park2023generative} demonstrate emergent social behaviors arising from agent interaction. Role assignment in multi-agent systems is a direct channel for bias: if persona or role assignment correlates with protected attributes, downstream behaviors of the assigned agent will differ systematically.

\textbf{4. Planning and task decomposition.} Agents decompose high-level goals into sequences of subtasks, often via chain-of-thought or structured planning procedures. The ordering and framing of subtasks can encode implicit assumptions about whose needs are primary, which resources are allocated first, and how ambiguous goals are resolved—all of which may vary systematically by demographic context.

\textbf{5. User modeling and personalization.} Agents adapt their behavior to individual users by inferring preferences, expertise, and context from interaction history or profile data. This adaptive capacity creates a personalization–stereotyping tension: inferences that improve average-case utility for a demographic group may simultaneously instantiate harmful stereotypes for individual users who deviate from that group's statistical profile.

\textbf{6. Long-horizon trajectory execution.} Over extended multi-turn interactions or autonomous task executions, small per-step biases can compound into large cumulative disparities. This longitudinal dimension is not captured by any existing single-turn benchmark and constitutes the most undertheorized component in current fairness research.

These six components—tool selection, memory and retrieval, multi-agent delegation, planning and decomposition, user modeling, and long-horizon execution—form the organizing axis of the taxonomy in Section 3. The central thesis of this survey is that fairness analysis must be conducted \emph{per component}, because the mechanism of harm, the appropriate measurement instrument, and the available mitigation strategy differ across components in ways that a global, post-hoc metric cannot capture.

\subsection{2.4 Scope and Methodology}

\textbf{Survey scope.} This survey covers empirical and theoretical research on unfairness and inequity that arises specifically within the architecture of LLM-based agent systems, understood as systems that perform multi-step, tool-augmented, or multi-agent reasoning toward a goal in a consequential domain. We do not survey the broader literature on bias in static LLM generation (single-turn QA, text classification, language modeling perplexity)—that literature is ably synthesized by Gallegos et al. \cite{gallegos2024bias} and Chu et al. \cite{chu2024fairness}—except insofar as those findings illuminate mechanisms operative in agentic settings. We also exclude non-LLM agents (classical planning systems, reinforcement learning agents, recommender systems without LLM backbones), though we draw on classical fair sequential decision-making results where they provide theoretical grounding.

\textbf{Inclusion criteria.} A paper was included if it met all of the following conditions: (i) the system studied uses an LLM as a primary reasoning component; (ii) the paper addresses bias, fairness, equity, or disparate treatment in outcomes or behaviors; (iii) the paper's findings bear on at least one of the six agent components identified above, or on the evaluation or mitigation of bias in such components. Papers exclusively studying static generation tasks without an agentic framing were excluded. Papers from non-LLM recommender systems or classical ML fairness were included only when directly cited by the agent fairness literature.

\textbf{Search and corpus construction.} The corpus was assembled through a systematic multi-branch search covering each of the six taxonomy components, plus evaluation and mitigation. Search queries were issued across arXiv (cs.AI, cs.CL, cs.LG), ACL Anthology, ACM DL, NeurIPS/ICML/ICLR proceedings, and FAccT proceedings, using variant terms for each component (e.g., "tool selection bias agent," "RAG fairness retrieval," "multi-agent bias delegation," "long-horizon fairness LLM"). Candidate papers were verified for existence, accurate metadata, and relevance before inclusion. The final corpus comprises \textbf{127 unique papers} spanning 2017–2026, with 61% published in 2024 or later, reflecting the field's rapid emergence. The corpus is organized across eight branches: foundations, LLM agents (general), tool/API selection bias, memory/retrieval bias, multi-agent delegation bias, planning/decomposition bias, user modeling and personalization harms, long-horizon fairness drift, evaluation metrics and benchmarks, and mitigation methods.

\textbf{Explicit out-of-scope items.} The following are excluded from analysis: (i) pure single-turn QA fairness benchmarks and their static LLM evaluations; (ii) non-LLM autonomous agents and classical planning fairness; (iii) fairness in LLM training data curation (addressed in the parent survey \cite{chu2024fairness}); (iv) adversarial robustness and safety unless the attack vector directly induces demographic disparities. Coverage of each in-scope paper across taxonomy components is visualized in the coverage matrix (Figure 2, Section 7).

\textbf{Positioning.} Three recent works partially overlap with our scope and must be distinguished carefully. Mohammadi et al. survey the evaluation and benchmarking of LLM agents broadly, treating fairness as one subdimension among many \cite{mohammadi2025evaluation}; our survey takes fairness as the \emph{sole} organizing axis and goes substantially deeper into each component's bias mechanisms. Vatsal et al. propose a seven-dimensional taxonomy for agentic AI in healthcare, in which fairness is one of seven dimensions and the treatment is largely conceptual rather than empirically grounded \cite{vatsal2026agentic}. Mayilvaghanan et al. apply counterfactual fairness evaluation to a contact-center agent in a single-domain deployment, introducing the CFR and MASD metrics \cite{mayilvaghanan2026counterfactual}; we generalize those metrics across all six agent components and multiple consequential domains. To our knowledge, no prior survey is dedicated to fairness organized by agent component across the full breadth of the LLM-based agent literature—a gap this work is designed to close.



\section{3. Where Bias Enters the Agent Pipeline}

The central claim of this survey is methodological: in an LLM \emph{agent}, fairness is not a property of a single forward pass but a property of a \emph{pipeline}. An agent perceives, selects tools, retrieves and accumulates memory, delegates subtasks across roles, plans and decomposes, models its user, and executes over long horizons. Each of these is a distinct locus where demographic disparity can originate, and each fails in a mechanistically different way. Organizing the literature by \emph{fairness metric} (group/individual/counterfactual) — as QA-era surveys do, including our parent \cite{chu2024fairness} — flattens this structure and obscures the fact that two agents with identical token-level bias scores can produce radically different \emph{action-level} disparities depending on where in the pipeline a bias is amplified or suppressed. We therefore organize the taxonomy by \textbf{agent component} (Axis 1). This is the dimension along which the three near-neighbor works do \emph{not} organize: \cite{mohammadi2025evaluation} treats fairness as one sub-topic of a broad agent-evaluation survey, \cite{vatsal2026agentic} as one of seven conceptual dimensions in a single (clinical) domain, and \cite{mayilvaghanan2026counterfactual} as a single-domain counterfactual method (contact-center QA) whose CFR/MASD machinery we generalize across \emph{all} components in §4.

This section defines six entry points (§§3.1–3.6) and then treats the cross-cutting question that distinguishes an agent from a chatbot: how component-level harms \textbf{compound} into an agentic multiplier (§3.7). Figure 1 renders these six points as a taxonomy tree over the agent anatomy of §2.3; Table 1 maps each component to its dominant harm mechanism and representative evidence. For each component we give a \emph{definition}, the \emph{mechanism of harm} (why the bias is component-specific, not merely inherited from the base model), the \emph{evidence} in the corpus, and the \emph{open edge} — the part of the component the field has not yet measured.

\subsection{3.1 Tool / API Selection Bias}

\textbf{Definition.} A tool-using agent must, at each step, choose \emph{which} external tool, API, retriever, or knowledge source to invoke. Tool-selection bias is systematic, demographically- or content-correlated skew in that choice — independent of, and often invisible to, the quality of the model's final text.

\textbf{Mechanism of harm.} Selection is a discrete decision over a tool registry, and LLMs make such selections with the same surface-form sensitivities that distort multiple-choice answering: option order, name salience, and description phrasing. \cite{pezeshkpour2023large} establishes the base pathology — LLM choice among enumerated options is highly sensitive to the \emph{order} in which options are presented, a positional bias that carries directly into tool registries presented as lists. \cite{blankenstein2025biasbusters} names and isolates the agentic form: across tool registries, LLMs exhibit a measurable \emph{tool-selection bias} toward particular tools regardless of fitness for the query, and the authors show it is both reproducible and mitigable, establishing that the disparity lives in the selection policy rather than the downstream answer. The harm is weaponizable: \cite{sneh2025tooltweak} demonstrates \emph{ToolTweak}, an attack that edits only a tool's surface metadata (name and description) to systematically inflate its selection rate, showing that whoever controls tool descriptions can steer an agent's routing — a fairness problem when tools correspond to providers, vendors, or sources with unequal demographic footprints. When the "tool" is a retriever, the bias becomes informational: \cite{dai2024bias} frames bias and unfairness in IR systems as a new class of LLM-era challenge, and \cite{wu2024rag} shows that the retrieval step itself — which corpus, which ranking — can inject unfairness even when the generator is held fixed. In a consequential clinical setting, \cite{xu2026ducx} (DUCX) decomposes unfairness in tool-using chest-X-ray agents and attributes a distinct share of the total disparity to the tool-invocation stage, separating it from model and data bias — the cleanest existing demonstration that tool selection is an \emph{independent} fairness locus.

\textbf{Open edge.} Existing work measures \emph{that} selection skews; almost none measures whether the skew correlates with the \emph{protected attributes of the user or subject} under a counterfactual swap. A tool-selection-disparity metric — does swapping a profile's demographics change which API the agent calls, holding the query fixed — does not yet exist as a benchmark, which is precisely the action-level gap §4.3 formalizes.

\subsection{3.2 Memory & Retrieval: Bias Accumulation Across Turns}

\textbf{Definition.} Agents persist state — conversation history, scratchpad memory, and retrieved documents — and condition each step on that accumulated context. Memory-and-retrieval bias is the introduction, \emph{accumulation}, or amplification of demographic skew through this conditioning, distinct from any bias in a single isolated response.

\textbf{Mechanism of harm.} Two sub-mechanisms operate. The first is \emph{retrieval injection}: \cite{hu2024free} ("No Free Lunch") shows that RAG can \emph{undermine} fairness even when the retrieved documents are individually innocuous and even for users who actively prompt for fairness, because the act of grounding shifts the generator's distribution; \cite{zhang2025evaluating} corroborates that retrieval augmentation measurably changes social-bias levels relative to the closed-book model; \cite{kim2025mitigating} localizes a controllable share of this to the \emph{embedder}, showing bias can be tuned at the retrieval-representation layer; and \cite{wang2025bias} demonstrates the adversarial extreme — poisoning the knowledge store to \emph{amplify} bias and steer downstream outputs. \cite{wu2024rag} again applies here as the systematic evaluation that RAG introduces unfairness. The second sub-mechanism is \emph{temporal accumulation}: bias is not merely present per-turn but \emph{compounds} as context grows. \cite{fan2024fairmt} (FairMT-Bench) is the anchor result — fairness degrades over \emph{multi-turn} dialogue in ways single-turn benchmarks miss, with models that pass single-turn tests failing as the conversation lengthens. \cite{geng2025accumulating} shows accumulating context literally \emph{changes the beliefs} of the model, shifting its stance as history grows; \cite{simhi2026old} characterizes the geometry — conversational history "geometrically traps" the model, making early framings self-reinforcing and hard to escape. These dynamics are exacerbated by positional pathologies: \cite{liu2023lost} ("Lost in the Middle") shows long-context models privilege information at the head and tail of the context window, so \emph{which} retrieved or remembered facts about a person get attended to is itself position-dependent and therefore skewable. Mitigation exists but is partial: \cite{kim2025mitigating} controls the embedder.

\textbf{Open edge.} The field measures bias accumulation in \emph{dialogue} (FairMT) but not yet in \emph{agentic memory} — the long-lived, cross-session stores that frame §3.6's drift. Whether bias compounds \emph{faster} under tool-augmented memory than under plain conversation is unmeasured.

\subsection{3.3 Multi-Agent Delegation & Role Assignment}

\textbf{Definition.} In multi-agent systems a controller assigns subtasks, personas, or roles to constituent agents and routes information among them. Delegation bias is disparity introduced by \emph{who is assigned what} — and by the social dynamics that emerge among role-bearing agents — over and above any single agent's bias.

\textbf{Mechanism of harm.} The seed is persona conditioning. \cite{gupta2023bias} ("Bias Runs Deep") shows that merely \emph{assigning a persona} to an LLM induces implicit reasoning biases that degrade performance on tasks unrelated to the persona — so role assignment is not neutral scaffolding but an active bias injection. \cite{cao2026from} traces this directly into the agentic regime, showing role assignment degrades agent \emph{robustness} and that biased chatbots become biased agents through their assigned roles. The harm then becomes \emph{systemic} through interaction. \cite{nguyen2025social} ("The Social Cost of Intelligence") is the central result: stereotypical bias \emph{emerges, propagates, and is amplified} across a multi-agent system, so the collective is more biased than its members — a genuinely agentic failure with no single-model analogue. \cite{madigan2025emergent} formalizes \emph{emergent} bias and fairness in multi-agent decision systems, showing fairness properties of the group are not predictable from the parts. \cite{li2025from} ("From Single to Societal") quantifies how persona-induced bias scales from one agent to the interacting society, and \cite{borah2024implicit} shows implicit bias surfaces specifically \emph{in the interaction} between agents, motivating detection and mitigation at the conversational seam rather than per-agent. Social structure makes it worse: \cite{campedelli2024want} shows LLMs in \emph{social-hierarchy} settings exhibit persuasion and anti-social behavior shaped by their rank, and \cite{mirza2025malibu} (MALIBU) provides a benchmark uncovering implicit bias specifically in multi-agent assignments. \cite{choi2025identity} closes the loop on mechanism by showing that \emph{anonymizing} agent identities in debate reduces bias — direct evidence that the disparity is caused by identity/role cues in the delegation, not by task content.

\textbf{Open edge.} Evidence is rich on \emph{whether} multi-agent setups amplify bias but thin on the \emph{delegation policy} as a controllable knob: there is no metric for "assignment disparity" (does a controller systematically route sensitive subtasks to lower-capability or more-biased sub-agents as a function of the subject's demographics), and the amplification factor of §3.7 has been demonstrated but not yet \emph{bounded}.

\subsection{3.4 Planning / Decomposition Disparities}

\textbf{Definition.} Agents plan: they decompose a goal into sub-goals, order steps, and reason explicitly (chain-of-thought, ReAct traces) before acting. Planning bias is demographic disparity that enters during this reasoning/decomposition stage — in the \emph{plan}, not merely the final answer.

\textbf{Mechanism of harm.} The counterintuitive finding is that reasoning does not cleanly remove bias and can \emph{introduce} it. \cite{wu2025reasoning} ("Does Reasoning Introduce Bias?") shows that adding explicit reasoning changes — and can worsen — social-bias outcomes relative to direct answering. \cite{zhou2025veracity} ("Veracity Bias") uncovers hidden beliefs that surface \emph{during} problem-solving reasoning: the model's intermediate steps encode demographic assumptions that a final-answer probe would not detect. The decisive reframing comes from \cite{li2025actions} ("Actions Speak Louder than Words"): when LLMs are made to \emph{decide and act} rather than answer, their agent decisions reveal implicit biases that text-level evaluations miss — direct evidence that the planning/acting stage is a distinct measurement surface and the empirical seed of this survey's measurement-gap thesis. \cite{parziale2026once} ("Once Upon a Team") instantiates the harm concretely: in LLM-driven software-team composition and task allocation, the planner allocates roles and tasks with demographic skew, a planning-stage disparity with direct hiring-adjacency. On the mitigation side, the same stage is a lever: \cite{kabra2025reasoning} shows reasoning-\emph{guided} fine-tuning can reduce bias, and \cite{hall2025guiding} guides decision-making with explicit \emph{fairness reward models} over the reasoning trace — both confirming that the plan is where intervention bites.

\textbf{Open edge.} No benchmark scores the \emph{plan itself} for disparity (e.g., are the sub-goals, ordering, or resource allocations a counterfactual function of the subject's demographics) independently of the final decision; planning disparity is observed anecdotally per-paper but not yet operationalized as an action-level metric.

\subsection{3.5 User Modeling & Personalization Harms}

\textbf{Definition.} Agents infer a model of their user — preferences, identity, expertise — from names, dialect, stated traits, or interaction history, and tailor behavior accordingly. Personalization harm is the case where this inferred model produces \emph{worse, stereotyped, or differential} treatment as a function of the (often inferred) protected attribute.

\textbf{Mechanism of harm.} The core tension is personalization-versus-stereotyping. \cite{kantharuban2024stereotype} ("Stereotype or Personalization?") shows that signaling user identity \emph{biases} chatbot recommendations toward stereotype rather than genuine preference — adaptation collapses into stereotyping. The trigger is often a mere \emph{name}: \cite{pawar2025presumed} ("Presumed Cultural Identity") shows names shape LLM responses by cuing a presumed cultural identity, and the privacy precondition is established by \cite{staab2023memorization}, which shows LLMs can \emph{infer} protected attributes (location, income, sex) from innocuous text — so an agent can model a user demographically without ever being told. Dialect is a second channel: \cite{lin2024assessing} demonstrates dialect \emph{unfairness} in reasoning tasks (performance drops on non-standard dialects), and \cite{mire2025rejected} ("Rejected Dialects") shows reward models penalize African American Language, baking the disparity into the alignment signal itself. The harm spans domains: \cite{weissburg2024llms} ("LLMs are Biased Teachers") finds personalized education delivers systematically different content quality by inferred student identity; \cite{deldjoo2024cfairllm} (CFaiRLLM) formalizes \emph{consumer} fairness in LLM recommenders, where personalization yields unequal recommendation utility across groups; and \cite{arzaghi2024understanding} surfaces intrinsic \emph{socioeconomic} bias that personalization can activate. Two papers complicate the picture usefully: \cite{vijjini2024exploring} exposes a safety-utility trade-off in personalized models (debiasing can cost helpfulness), and \cite{tonneau2026different} ("Different Demographic Cues Yield Inconsistent Conclusions") is a crucial methodological caution — \emph{which} demographic cue you use (name vs. explicit label vs. dialect) changes the measured personalization bias, so single-cue audits can mislead.

\textbf{Open edge.} \cite{tonneau2026different} implies the field's personalization findings are cue-fragile, yet no agentic benchmark varies the cue \emph{channel} systematically; and personalization harm has been measured in single responses, not over the persistent user model an agent maintains across a session — the bridge to §3.6.

\subsection{3.6 Long-Horizon Fairness Drift Over Trajectories}

\textbf{Definition.} Agents operate over extended trajectories and, when deployed, over feedback loops in which their own outputs become future inputs. Long-horizon fairness drift is the \emph{temporal} degradation — or self-amplification — of fairness across turns, sessions, or deployment cycles, even when the per-step model is fixed.

\textbf{Mechanism of harm.} Three loops drive drift. First, \emph{intra-trajectory} drift within a single long interaction: \cite{fan2024fairmt} again anchors the multi-turn degradation, and \cite{ma2026implicit} shows implicit bias \emph{accumulates and propagates} specifically in LLM long-term memory, making the agent more biased the longer it persists state. Sycophancy is a named driver: \cite{liu2025truth} ("Truth Decay") quantifies multi-turn sycophancy — the model progressively conforms to the user — and \cite{hong2025measuring} measures sycophancy growth across multi-turn dialogue, a mechanism by which a user's biased framing is incrementally adopted by the agent. Second, \emph{deployment feedback loops}: \cite{wyllie2024fairness} ("Fairness Feedback Loops") shows training on a model's own synthetic outputs \emph{amplifies} bias across generations; \cite{wang2026observations} characterizes this self-consuming \emph{performative} loop for LLMs specifically and offers remedies; and \cite{mansoury2020feedback}, from the recommender-systems lineage, gives the canonical feedback-loop bias-amplification result that the LLM work now inherits. Third, the \emph{sequential-decision} foundations that name the formal target: \cite{hu2022achieving} (Hu and Zhang) defines and achieves \emph{long-term} fairness in sequential decision-making; \cite{xu2024adapting} adapts static fairness notions to the sequential setting via an \emph{equal long-term benefit rate}, showing that single-step fairness constraints can \emph{fail} to deliver fairness over a trajectory; and \cite{she2025fairsense} (FairSense) provides simulation-based \emph{long-term} fairness analysis of ML-enabled systems, the systems-level tooling an agentic audit would need. \cite{liu2025truth} and \cite{hong2025measuring} together establish that drift is measurable per-turn.

\textbf{Open edge.} The sequential-fairness theory (\cite{hu2022achieving}, \cite{xu2024adapting}) and the LLM-drift empirics (\cite{fan2024fairmt}, \cite{ma2026implicit}) have not been joined: no work yet evaluates an \emph{LLM agent} against a formal long-term-fairness criterion over a realistic multi-session deployment, leaving the field's most consequential fairness regime essentially unmeasured.

\subsection{3.7 Cross-Cutting: How Component Harms Compound — The Agentic Multiplier}

The six entry points are not additive. The defining feature of an \emph{agent} — and the reason QA-level fairness scores are insufficient — is that a bias originating in one component is \emph{propagated and amplified} by the next, so total trajectory-level disparity exceeds the sum of per-component disparities. We call this the \textbf{agentic multiplier}, and two corpus results make it concrete rather than rhetorical. \cite{nguyen2025social} demonstrates the multiplier directly within the multi-agent component: stereotypical bias \emph{emerges} where none was salient, \emph{propagates} between agents, and is \emph{amplified} by their interaction — a strictly super-additive outcome relative to the constituent models. \cite{madigan2025emergent} generalizes the point: fairness in multi-agent decision systems is \emph{emergent}, not inherited, so component-level audits cannot be composed into a system-level guarantee.

The same logic links the \emph{other} components pairwise. A demographic attribute \emph{inferred} by user modeling (§3.5, via \cite{staab2023memorization}) becomes a \emph{retrieved and accumulated} fact in memory (§3.2, \cite{fan2024fairmt}, \cite{ma2026implicit}); it then conditions \emph{tool selection} (§3.1) and \emph{plan decomposition} (§3.4, \cite{li2025actions}); a multi-agent controller \emph{delegates} the resulting subtask along role lines (§3.3, \cite{nguyen2025social}); and the whole loop \emph{drifts} over the trajectory (§3.6, \cite{wyllie2024fairness}, \cite{xu2024adapting}). Each hand-off is an opportunity for amplification, and — critically — a token-level fairness probe at any single hand-off can read as clean while the \emph{action} the agent ultimately takes is sharply disparate, exactly the failure \cite{li2025actions} exposes when it moves evaluation from words to decisions.

This compounding is the structural reason the survey is organized by component rather than by metric, and it is the conceptual core of the measurement gap developed in §4: today's benchmarks score answers at one point in the pipeline, but the harm is a property of the \emph{trajectory} across all six. Table 1 summarizes each component's mechanism and evidence; Figure 1 situates them on the agent anatomy; and the coverage analysis of §7 (Figure 2) shows that the cells most relevant to the agentic multiplier — action-level, counterfactual, long-horizon — are precisely where the field is emptiest. No prior survey, including our parent \cite{chu2024fairness} and the three near-neighbors \cite{mohammadi2025evaluation}, \cite{vatsal2026agentic}, and \cite{mayilvaghanan2026counterfactual}, organizes the literature around this compounding pipeline; doing so is this survey's primary contribution.



\section{4. How to Evaluate Agent Fairness}

The translation from LLM-level fairness metrics to agent-level evaluation is neither direct nor complete. Single-turn benchmarks measure what a model \emph{says}; agent pipelines must be evaluated on what a system \emph{does} — which tools it selects, whose requests it escalates, how disparities compound across turns. This section maps existing metric families onto agentic settings (§4.1), surveys counterfactual evaluation methods now being adapted for agent trajectories (§4.2), catalogs the emerging but incomplete literature on action-level disparity (§4.3), and takes inventory of current datasets and benchmarks, documenting the central measurement gap: nearly every available resource operates at the QA level (§4.4, Table~2).

\subsection{4.1 Metric Families Mapped onto Agent Settings}

Three classical fairness notions — group, individual, and counterfactual — each impose distinct demands when lifted from single-turn scoring to agent trajectories.

\textbf{Group fairness} in the LLM literature is typically operationalized as demographic parity or equalized odds across protected attributes, measured over a held-out evaluation set of prompts \cite{hardt2016equality,mehrabi2021survey}. At the agent level, the relevant outcome is not a token distribution but a decision distribution: approval vs. denial rates across loan applicants \cite{azime2025accept}, escalation rates to a human reviewer across patient groups \cite{adappanavar2025mfarm}, or tool-invocation frequencies across demographic conditions. The group changes; the question of whether outcome distributions are equalized remains, but the outcome itself is now an action rather than a prediction. Surveys of LLM-level fairness measurement \cite{gallegos2024bias,chu2024fairness} acknowledge group disparity metrics but are organized around textual outputs, not action logs.

\textbf{Individual fairness} demands that similarly situated inputs receive similar treatment \cite{dwork2012fairness}. In an agentic context, individual fairness applies across entire interaction trajectories: two agents processing otherwise-identical applicant profiles that differ only in a protected attribute should reach the same decision along the same decision path. This trajectory-level generalization is rarely formalized in current evaluation practice. The closest instantiation in the corpus is the consistency-based framing of \cite{basu2026names}, which measures whether intervention recommendations change when names (and thus inferred ethnicity and gender) are altered while holding clinical features constant, finding systematic inconsistency rates that aggregate statistics would mask.

\textbf{Counterfactual fairness} — no individual should receive a different outcome in a world where only their protected attribute differed \cite{kusner2017counterfactual} — translates most naturally to agents, because the demographic swap can be applied to the full input to an agent and the full output trajectory observed. This generalization is now actively pursued (§4.2), and it is the evaluation paradigm most capable of surfacing component-level disparity.

A fourth family, \textbf{reward-model and alignment-level group fairness}, is surfacing as agents are increasingly evaluated or steered via reward models \cite{song2025large}. If the reward model itself encodes demographic preferences — as \cite{song2025large} documents across eight groups — then agent behaviors shaped by that signal will inherit the bias. Evaluation must therefore propagate upstream from the model's action into the mechanism that rewards it.

The critical gap across all three families is the \emph{action granularity} of the denominator: existing frameworks aggregate across many prompts to compute a disparity statistic but treat each agent response as atomic, ignoring which internal component produced the disparity-generating behavior. An agent can return an equalized final answer through biased tool routing; group-fairness statistics at the response level will not detect this.

\subsection{4.2 Counterfactual Methods for Agents}

Counterfactual evaluation has seen the most rapid methodological development at the agent level, producing several distinct frameworks in the 2025–2026 period.

\textbf{CFR and MASD} (Counterfactual Flip Rate and Mean Absolute Score Difference) are introduced by \cite{mayilvaghanan2026counterfactual} for evaluating an LLM-based quality-assurance agent in a contact-center deployment. CFR counts the fraction of agent evaluations that reverse when the agent's input is perturbed by a demographic swap alone; MASD measures the mean absolute difference in continuous quality scores across paired counterfactual conditions. These two statistics — one binary, one continuous — cleanly decompose the fairness question: is the \emph{decision} sensitive to protected attributes, and by how much? The present survey generalizes both metrics from the customer-service domain to any agent component: a tool-selection agent can be probed with CFR over tool-choice outcomes; a planning agent can be probed over plan step count or resource allocation. \cite{mayilvaghanan2026counterfactual} demonstrate non-trivial CFR values in their deployment setting, motivating this generalization.

\textbf{CAFFE} (Counterfactual Fairness Evaluation) \cite{parziale2025systematic} provides a framework for systematic counterfactual generation and evaluation across a broad set of LLM behaviors, going beyond single demographic attributes to structured, multi-attribute perturbations. The companion \textbf{SCOPE} dataset \cite{parziale2026scope} supplies stereotyped prompts specifically designed to stress-test counterfactual fairness, including prompts covering protected attributes that existing benchmarks leave sparse. Together, CAFFE and SCOPE constitute a test harness that can be applied to agent responses at any pipeline stage where a natural-language interface is exposed.

\textbf{Flip-rate auditing} \cite{basu2026names} operationalizes the same core intuition in a high-stakes clinical and judicial context: names carrying demographic signals are substituted in otherwise-identical decision scenarios, and the fraction of cases in which the agent's recommendation flips is recorded. \cite{basu2026names} find that name-based demographic signals alter verdicts and intervention recommendations at rates that expose systematic bias invisible to aggregate-group statistics. Because flip-rate probing requires only a demographic-perturbation function and access to the agent's action-level output, it generalizes across agent architectures.

\textbf{FiSCo} (Fairness in Semantics and Counts) \cite{xu2025quantifying} reframes fairness measurement beyond the token level, operating on semantic units and statistical summaries of agent outputs rather than raw text. This is important for multi-step agent pipelines in which the final answer may be verbally equalized while the reasoning trace and intermediate actions carry demographic differentiation. FiSCo provides a mechanism to detect this form of surface-level parity masking underlying disparity.

\textbf{Equal-Access auditing} \cite{amirimargavi2026equal} conducts a counterfactual audit across a set of LLM deployments, examining whether access quality — the comprehensiveness and actionability of responses — differs by inferred user demographics. In agentic settings, this maps directly onto service-quality disparity: whether the agent retrieves the same depth of information, invokes the same complement of tools, or routes the request with the same priority across demographic groups.

Collectively, these frameworks share a common scaffold: (1) construct demographically paired input sets, (2) observe agent outputs or trajectories under each condition, (3) compute a distance or flip statistic. What they do not yet provide is systematic coverage of \emph{intermediate} agent states — the tool selected at step 3, the memory retrieved at step 5, the subagent delegated to at step 2. Counterfactual evaluation of agent trajectories remains largely endpoint-to-endpoint; the component-level localization of bias is a standing open problem addressed in §8.

\subsection{4.3 Action-Level Disparity Metrics}

A small but growing body of work insists that agent fairness must be measured at the level of \emph{actions taken}, not answers produced. Two papers in the corpus develop explicit action-level metrics.

\cite{li2025actions} operationalizes the gap directly in its title and content: agent decisions — the actual choices an agent makes in structured decision tasks — reveal implicit biases that are invisible at the verbal-output level. The paper demonstrates that models presenting verbally neutral or explicitly equalized language simultaneously issue actions (structured recommendations, selections, rankings) that carry demographic disparity signals, providing empirical evidence that evaluating textual outputs alone understates actual agent-level bias.

\cite{blankenstein2025biasbusters} introduces the construct of \emph{tool selection bias}: systematic differences in which tools an LLM-based agent selects as a function of the demographic framing of the input query. The paper demonstrates that tool-selection patterns — not the downstream answers — are the proximate locus of the disparity, and proposes metrics that track tool-invocation frequencies across demographic conditions rather than aggregating over the final response. This operationalizes what we term \emph{decision-path disparity} in the taxonomy of §3.1.

What is currently missing is a standardized action-level disparity metric vocabulary. The field needs at minimum three metrics not yet systematically defined: (i) \textbf{tool-invocation disparity} — whether the distribution over tool choices differs by demographic group; (ii) \textbf{escalation-rate disparity} — whether requests are routed to a human reviewer or higher-capability resource at different rates across groups; and (iii) \textbf{plan-length and resource-allocation disparity} — whether planning agents allocate different step counts, tool calls, or downstream resources to otherwise equivalent requests from different demographic groups. None of these has a canonical formula, a reference implementation, or a benchmark split designed to surface it. This is the sharpest edge of the measurement gap.

\subsection{4.4 Datasets and Benchmarks}

Table~2 provides an inventory of fairness datasets and benchmarks in the corpus, organized by fairness level (intrinsic/extrinsic), domain, and whether the benchmark evaluates agentic behavior. The punchline is immediate: with very few partial exceptions, every resource evaluates QA-level behavior, not agent actions.

\textbf{Clinical and biomedical benchmarks.} \cite{xiao2025fairmedqa} (FairMedQA) benchmarks bias in LLMs applied to medical question answering across demographic subgroups, providing a well-structured evaluation set for a consequential domain. \cite{pfohl2024toolbox} (published in \emph{Nature Medicine}) provides a broader toolbox for surfacing health equity harms, covering both intrinsic model biases and extrinsic task performance disparities in clinical decision settings. \cite{ghosh2025medequalqa} (MedEqualQA) introduces counterfactual reasoning into the medical QA fairness paradigm, pairing questions with demographic perturbations to measure consistency. \cite{adappanavar2025mfarm} (mFARM) extends fairness evaluation to multi-faceted harm assessment in clinical decision support, covering intersecting demographic attributes and multiple harm categories. \cite{wu2024fmbench} (FMBench) benchmarks fairness in multimodal LLMs on medical tasks, adding a vision modality that is increasingly relevant as agent pipelines incorporate imaging tools. None of these benchmarks includes an agent-execution layer: they probe model responses to clinical questions, not agent decisions within clinical workflows.

\textbf{Decision-domain benchmarks.} \cite{azime2025accept} evaluates LLM fairness in loan approval decisions across table-to-text serialization approaches, examining how input representation affects demographic disparity in lending recommendations. \cite{hu2025llms} (LLMs on Trial) evaluates judicial fairness in LLM-generated assessments, documenting differential treatment across defendant demographics in structured sentencing and evaluation scenarios. These benchmarks are closer to the agentic use case — they involve consequential decisions — but still elicit single-turn model responses rather than multi-step agent trajectories.

\textbf{General-purpose fairness evaluation resources.} \cite{zhang2025datasets} surveys the landscape of fairness datasets for language models, organizing them along dimensions of protected attribute, task type, and measurement approach — an essential reference for benchmark selection. \cite{wang2024ceb} (CEB) provides a compositional evaluation benchmark that combines multiple fairness dimensions (stereotype, toxicity, disparate performance) into structured evaluation suites. \cite{fan2024fairmt} (FairMT-Bench) focuses specifically on multi-turn dialogue fairness, making it the closest existing resource to an agentic evaluation benchmark: it examines whether fairness degrades across conversation turns, which is a restricted analogue of long-horizon fairness drift (§3.6). \cite{song2025large} benchmarks group fairness specifically in reward models, targeting the alignment mechanism rather than the deployed model. \cite{bouchard2024bring} (LangFair) provides a use-case-specific bias and fairness evaluation toolkit that can be applied to production LLM deployments, supporting customizable evaluation pipelines.

\textbf{The measurement gap.} The inventory in Table~2 reveals a consistent structural absence. Every benchmark in the clinical QA tier probes a model's answer to a clinical question. Every decision-domain benchmark elicits a model's rating or recommendation in a single prompt. The multi-turn resources (FairMT-Bench) cover conversational dynamics but not tool use, memory retrieval, or delegation. No benchmark in the corpus — by design — presents a fairness evaluation problem at the level of a full agent trajectory: which tool was selected, what was retrieved from memory, how was the task decomposed, which subagent was assigned which role, and did these intermediate actions differ by demographic group? The measurement infrastructure is one abstraction layer below where the field needs it to be. This gap is not merely an absence of data; it reflects a disciplinary boundary between the NLP fairness tradition (which measures outputs) and the systems tradition (which measures behaviors). Closing it requires benchmarks that expose and instrument agent internals, not just agent answers. This agenda is taken up in §8.1, and our empirical audit in §6 provides a direct demonstration that QA-level evaluation systematically fails to surface action-level disparity. The companion FairMedAgent benchmark is designed specifically to address this gap.


\noindent\textbf{Table~2 — Benchmark Inventory} is structured with columns: \textit{Benchmark}, \textit{Fairness Level} (intrinsic/extrinsic), \textit{Domain}, \textit{Agentic?} (yes/no/partial), \textit{Attributes Covered}, and \textit{Key Citation}. Representative entries: FairMedQA (extrinsic, clinical, no, race/gender/age, \cite{xiao2025fairmedqa}); mFARM (extrinsic, clinical, partial, multi-attribute, \cite{adappanavar2025mfarm}); MedEqualQA (extrinsic, clinical, no, counterfactual, \cite{ghosh2025medequalqa}); FMBench (extrinsic, medical-multimodal, no, race/gender, \cite{wu2024fmbench}); Pfohl Toolbox (extrinsic, clinical, no, race/gender/SES, \cite{pfohl2024toolbox}); Accept-or-Deny (extrinsic, lending, no, race/gender, \cite{azime2025accept}); LLMs on Trial (extrinsic, judicial, no, race/gender, \cite{hu2025llms}); CEB (intrinsic+extrinsic, general, no, multi-attribute, \cite{wang2024ceb}); FairMT-Bench (extrinsic, dialogue, partial, multi-attribute, \cite{fan2024fairmt}); SCOPE (counterfactual, general, no, multi-attribute, \cite{parziale2026scope}); LangFair (extrinsic, general, no, configurable, \cite{bouchard2024bring}); Reward-model benchmark (intrinsic, general, no, group, \cite{song2025large}). No row in the table bears "yes" in the Agentic column — this is the figure that the section is designed to make unmistakable.



\section{5. How to Mitigate}

The mitigation landscape for agent fairness is expansive on paper but narrow in practice: the overwhelming majority of proposed interventions target the underlying language model at the token or decoding level, not the agentic components — tool selectors, memory modules, planners, or delegation protocols — where the taxonomic harms documented in Section 3 actually originate. Table 3 maps each mitigation family to the agent component it touches, the pipeline stage at which it operates, and the evidence supporting it. The central takeaway from that mapping is that the agent layer remains largely unaddressed.

\subsection{5.1 Prompting and In-Context Debiasing}

The most accessible and deployment-friendly family of mitigations modifies the prompt itself, inserting fairness-sensitive instructions without altering model weights. \cite{li2024prompting} frames this through a causal lens: by integrating front-door adjustment logic into the prompt construction process, the method guides the LLM to generate outputs that are less confounded by protected-attribute signals already present in the input context. Complementing this, \cite{zhang2024causal} formalizes causal prompting as a systematic application of front-door adjustment at inference time, demonstrating that structuring the chain of reasoning to block spurious demographic pathways reduces downstream bias in open-ended generation tasks. A different structural approach comes from \cite{furniturewala2024thinking}, whose structured prompting framework ("Thinking Fair and Slow") finds that multi-step, slow-deliberation templates — inspired by dual-process cognitive theory — consistently outperform flat zero-shot fairness instructions across demographic categories. \cite{du2024causal} extends the causal paradigm into active learning: rather than prompting alone, the method selects which labeled examples to include in-context based on causal influence scores, steering the model away from spurious demographic correlations in few-shot settings.

From an agent-component perspective, prompting interventions are most naturally applied to the planning and reasoning stage (§3.4) and to user-facing response generation (§3.5). However, their reach is structurally limited: they cannot directly constrain which tools an agent selects (§3.1), what a retrieval index surfaces (§3.2), or how delegation authority is distributed in a multi-agent pipeline (§3.3). Any agent that overwrites or reformats the system prompt between turns further degrades prompt-level guarantees. Prompting mitigations are therefore best understood as a lightweight first layer, not a complete solution.

\subsection{5.2 Fine-Tuning and Alignment}

Training-time mitigations offer deeper integration than prompting by shaping model weights directly. \cite{ouyang2022training} introduced RLHF as the canonical alignment paradigm, using human preference data to steer outputs toward helpfulness and harmlessness; while not fairness-specific by design, RLHF is the backbone on which most subsequent bias-aware alignment work is constructed. A distributional critique emerges from \cite{xiao2024algorithmic}, which formally characterizes RLHF's vulnerability to preference collapse — a failure mode in which majority-group preferences dominate the reward signal and model outputs systematically diverge from minority preferences. \cite{chakraborty2024maxmin} addresses this directly through MaxMin-RLHF, replacing the standard reward-maximization objective with a maximin criterion that optimizes for the worst-off preference subgroup, producing alignment that is robust across diverse human populations rather than calibrated to the plurality.

Direct Preference Optimization (DPO) has emerged as a computationally cheaper alternative to RLHF; \cite{allam2024biasdpo} adapts DPO explicitly for bias reduction, constructing preference pairs from contrastive biased/debiased completions and fine-tuning to prefer the debiased response. \cite{kabra2025reasoning} takes a complementary approach: rather than optimizing a fairness objective post-hoc, the method incorporates fairness-aware chain-of-thought supervision during fine-tuning, so that the model learns to reason through demographic parity considerations before generating a decision. This reasoning-guided approach is particularly relevant for planning tasks (§3.4), where an agent decomposes a problem into sub-tasks that may receive unequal treatment across demographic groups.

These alignment methods carry meaningful limitations for agent deployment. Fine-tuning and RLHF operate on the base language model; once that model is embedded in a multi-agent framework or given tool access, the alignment guarantees are not automatically transferred to the agent's action selection or delegation logic. Aggregate reward signals do not distinguish between a bias introduced at the planning step and one introduced at the tool-call step, making component-level attribution — and targeted correction — effectively invisible at training time.

\subsection{5.3 Guardrails and Filters}

Inference-time guardrails interpose a protective layer between the model's raw output and the downstream consumer. \cite{dong2024building} provides a systematic taxonomy of LLM guardrail architectures, distinguishing input filters (blocking demographic-sensitive queries), output filters (post-hoc screening of model completions), and architectural wrappers (modular safety classifiers surrounding the core model). For agent systems, the relevant extension is guardrails applied not to conversational output but to action proposals: before an agent executes a tool call or delegates a sub-task, a classifier audits the proposed action for demographic disparity. In principle, this approach is component-agnostic — it can be inserted at tool selection (§3.1), at delegation boundaries (§3.3), or at long-horizon checkpoints (§3.6).

In practice, guardrail deployments in agent contexts remain immature. Existing systems largely re-use classifiers trained on natural language toxicity or refusal tasks; none of the corpus papers report a guardrail specifically trained on action-level fairness signals. The measurement gap identified in Section 4 is directly implicated: without action-level fairness datasets, training action-level guardrails is not tractable.

\subsection{5.4 Multi-Agent Debiasing}

Multi-agent architectures introduce new bias pathways (§3.3), but they also create structural opportunities for peer-level correction that single-agent systems cannot exploit. \cite{owens2024multi} proposes a general multi-LLM debiasing framework in which multiple agents generate responses independently and a synthesis step aggregates across models, suppressing idiosyncratic biases through ensemble diversity. \cite{xu2024mitigating} frames this as a multi-objective optimization problem within a multi-agent setting: specialized agents are assigned to maximize different fairness criteria, and a coordinator agent reconciles their proposals, achieving bias reduction without sacrificing task utility.

Two papers address cultural and identity-linked biases specifically. \cite{ki2025multiple} organizes LLM agents into a structured debate protocol oriented toward equitable cultural alignment, finding that adversarial multi-agent argumentation surfaces and corrects cultural stereotypes more effectively than single-agent self-critique. \cite{choi2025identity} diagnoses a distinct failure mode: when agents are assigned or infer social identities, those identities systematically skew the direction of multi-agent reasoning in ways consistent with demographic bias. The proposed mitigation — selective anonymization of agent identity cues within the debate context — reduces identity-driven skew without eliminating the epistemic benefits of role differentiation.

At the implicit level, \cite{borah2024implicit} demonstrates that implicit biases circulate through multi-agent interaction histories and proposes a detection-and-mitigation pipeline that analyzes interaction logs to identify bias propagation patterns before they crystallize into persistent behavioral tendencies. Taken together, these approaches occupy a unique position in Table 3: they are the only mitigation family that operates natively at the agent-architecture level rather than at the LLM token level. Their limitation is symmetrical with their strength — they require a multi-agent infrastructure to function and are not portable to single-agent deployments.

\subsection{5.5 Process-Level and Architectural Interventions}

The most architecturally integrated category of mitigations embeds fairness enforcement into the agent's operational logic rather than applying it as an external correction. This subsection covers three distinct loci: inference-time activation-level steering, RAG-level retrieval fairness, and memory-level bias hygiene.

\textbf{Inference-time steering.} \cite{bai2022constitutional} introduced Constitutional AI, in which a model is trained to self-critique and revise its outputs against a fixed set of normative principles ("the constitution"); the procedural discipline of self-revision produces systematically less harmful and more fair outputs than standard RLHF. \cite{huang2024collective} extends this paradigm by replacing a top-down authored constitution with one derived from public deliberation, using collective human input to define the normative constraints — a meaningful shift for fairness because the resulting constitution is less likely to encode majority-group values by default. At the inference stage, \cite{li2025fairsteer} operationalizes activation-level debiasing: FairSteer uses dynamic activation steering — identifying fairness-relevant directions in the model's representation space and suppressing activations along those directions at inference time — achieving debiasing without weight updates. This approach is compatible with agent deployment because it operates on the forward pass, not on training, and can in principle be conditioned on agent-state signals.

\textbf{RAG-level interventions.} The retrieval component is a high-priority bias entry point (§3.2). \cite{kim2024fair} provides the first systematic analysis of how ranking fairness in the retrieval stage propagates into downstream generation quality across demographic groups, establishing that retrieved-document demographic representation directly shapes output fairness. \cite{kim2025mitigating} translates this insight into a practical intervention: by intervening in the embedding model rather than the ranker or generator, the method corrects for demographic skew at the source, reducing bias without requiring post-hoc re-ranking or generation-level filtering.

\textbf{Memory-level hygiene.} Long-term agent memory is among the least-studied bias propagation channels. \cite{ma2026implicit} characterizes how implicit biases accumulate and compound across LLM long-term memory — a model of bias accumulation that is structurally analogous to feedback amplification in recommender systems \cite{mansoury2020feedback} — and proposes memory audit and selective forgetting mechanisms as remediation. The memory hygiene approach is unique in that it addresses §3.6 (long-horizon fairness drift) directly, targeting the precise mechanism by which per-turn biases compound over trajectories.

\textbf{Summary assessment.} The process-level tier is the most promising direction for genuinely agent-level fairness intervention, but it is also the most sparsely populated in the corpus. As Table 3 makes clear, the preponderance of mitigation evidence clusters at the LLM/token level (prompting, alignment, guardrails), with multi-agent debiasing and architectural interventions representing a small and recent minority. A standing gap is the absence of any mitigation that jointly addresses tool selection bias (§3.1), memory accumulation (§3.2), and long-horizon drift (§3.6) within a unified architectural framework. Building such a framework — one that treats fairness as an agent-level property enforced across the full component stack — is the central open problem addressed in Section 8.6.



\section{6. Original Empirical Study: A Cross-Framework Bias Audit}

The preceding sections argue, from the literature, that fairness in LLM-based agents is a property of \emph{components and actions}, not of single-turn answers. That argument is only as strong as its empirical purchase. To furnish a concrete, reproducible demonstration — and to satisfy the requirement that a survey carry original results — we register a compact \textbf{cross-framework bias audit}: the same consequential-decision task is routed through agent configurations of increasing architectural complexity, and disparity is measured not only at the level of the final verdict (as QA-style audits do) but at the level of the \emph{intermediate actions} the agent takes. The study is deliberately compute-light (pinned models, a single estimated API cost of roughly \$100--300, synthetic profiles), so that its value is methodological rather than scale-dependent: it shows that component- and action-level disparity is \emph{measurable} and \emph{present} where answer-level scoring is silent. The full protocol is fixed below; the run is pending, and numerical results are explicitly deferred to camera-ready (Section 6.5).

\subsection{6.1 Experimental Setup}

\textbf{Task and profiles.} We instantiate a single consequential-decision task — resume screening for a fixed job requisition (with loan-application triage as a pre-registered replication domain) — chosen because answer-level bias here is already well documented for non-agentic LLMs and thus provides a calibrated baseline against which agentic effects can be isolated. Name- and demographic-based disparities in LLM hiring recommendations are established by \cite{nghiem2024you}, \cite{veldanda2023emily}, \cite{wilson2024gender}, and \cite{an2025measuring}; the analogous evidence for lending and credit is given by \cite{bowen2024measuring}, \cite{feng2023empowering}, and \cite{azime2025accept}. Building on this base lets us attribute any \emph{additional} disparity to the agent scaffolding rather than to the underlying model.

To avoid any human-subjects exposure, all candidate profiles are \textbf{synthetic} (no real applicants, no IRB review required). We construct a small set of base profiles holding qualifications fixed, then generate \textbf{counterfactual variants} that differ only in demographic signal — primarily via names as a proxy for race and gender, following the name-perturbation designs of \cite{nghiem2024you} and \cite{pawar2025presumed}, and secondarily via explicit demographic cues. Because demographic-cue \emph{encoding} is itself a confound — \cite{tonneau2026different} show that different cue formats (names vs. self-disclosure vs. inferred attributes) yield inconsistent bias conclusions — we vary the cue channel as a controlled factor rather than fixing a single representation.

\textbf{Agent configurations.} The independent variable is \emph{agent architecture}. We compare four configurations of escalating compositional complexity, holding the base model, prompt scaffold, and decision rubric constant:

1. \textbf{C0 — Single-call LLM (no agency).} The base model emits a verdict directly. This is the QA-level baseline against which all agentic deltas are measured.
2. \textbf{C1 — Single-LLM tool-use.} One agent equipped with tools (e.g., a skills/keyword matcher, a policy-lookup tool, a scoring calculator), selecting among them per the function-calling paradigm of \cite{schick2023toolformer}. This isolates \textbf{tool-selection} as a bias entry point, the mechanism characterized by \cite{blankenstein2025biasbusters} and \cite{li2025actions}.
3. \textbf{C2 — ReAct-style reasoning-and-acting agent.} Interleaved reasoning traces and tool calls over multiple turns \cite{yao2023react}, adding \textbf{planning/decomposition} and accumulated-context channels; the latter is the locus of the multi-turn drift documented by \cite{fan2024fairmt} and \cite{geng2025accumulating}.
4. \textbf{C3 — Two-agent delegation/debate.} A delegation or debate configuration in which roles are assigned across two agents \cite{wu2023autogen} (e.g., a screener and an adjudicating critic). This exposes \textbf{multi-agent delegation and role-assignment} effects, the channel implicated by \cite{madigan2025emergent}, \cite{nguyen2025social}, and \cite{cao2026from}.

Each configuration processes the identical profile set; models are \textbf{version-pinned} and decoding parameters fixed across conditions so that observed differences are attributable to architecture, not to model drift or sampling. The design is a within-profile, between-architecture comparison, enabling paired counterfactual contrasts.

\subsection{6.2 Metrics}

We measure disparity at two levels, deliberately separating the \emph{answer} layer (where current benchmarks operate) from the \emph{action} layer (where the survey argues the gap lies).

\textbf{Answer-level (verdict) disparity.} We adopt two counterfactual measures generalized from the contact-center auditing framework of \cite{mayilvaghanan2026counterfactual}: the \textbf{Counterfactual Flip Rate (CFR)} — the fraction of base profiles whose final verdict changes when only the demographic signal is swapped — and the \textbf{Mean Absolute Score Difference (MASD)} — the mean absolute change in the agent's continuous score (e.g., recommendation confidence or rank) across counterfactual pairs. Reporting both a discrete-flip and a magnitude metric follows the intervention-consistency logic of \cite{basu2026names}.

\textbf{Action-level disparity.} The contribution of this audit is to compute disparity over the agent's \emph{intermediate behavior}, not only its output. For each configuration that takes actions (C1--C3), we record, per counterfactual pair, the divergence in: (i) \textbf{tool-selection rate} — which tools are invoked and how often, the disparity surfaced for tool-using agents by \cite{blankenstein2025biasbusters} and, in the clinical-imaging setting, decomposed by \cite{xu2026ducx}; (ii) \textbf{decision-action rate} — disparities in the operational action taken (advance/reject, approve/deny/refer), per the agent-action framing of \cite{li2025actions}; and (iii) \textbf{escalation/deferral rate} — how often the agent defers, requests more information, or escalates to a human, a margin where differential treatment can hide behind nominally identical final verdicts. These action-level quantities are computed as counterfactual-pair disparities directly parallel to CFR/MASD, so that answer- and action-level disparity are reported on a comparable footing within a single table.

\subsection{6.3 Hypotheses}

The protocol is registered against two falsifiable hypotheses, each tied to a mechanism in the taxonomy of Section 3.
\item \textbf{H1 (the measurement gap is real).} There exist configurations in which answer-level disparity (CFR/MASD at the final verdict) is small or statistically indistinguishable from the C0 baseline, while action-level disparity (tool, decision, or escalation rates) is materially larger. In other words, component-level disparity is \emph{invisible to QA-level scoring}. Confirmation would empirically substantiate the survey's central claim that scoring answers under-measures agent unfairness; the action-as-signal premise is drawn from \cite{li2025actions}.
\item \textbf{H2 (multi-agent composition compounds bias).} Disparity is non-decreasing in architectural complexity along the ordering C0 ⪯ C1 ⪯ C2 ⪯ C3, with the two-agent delegation/debate configuration (C3) exhibiting the largest disparity. This operationalizes the \emph{agentic multiplier} of Section 3.7: the emergence and amplification of bias across interacting agents reported by \cite{madigan2025emergent} and \cite{nguyen2025social}, and the role-assignment sensitivity documented by \cite{cao2026from} and \cite{li2025from}. As a registered counter-hypothesis, multi-agent \emph{debate} can also reduce disparity when identity cues are anonymized \cite{choi2025identity} or when debate is used as a debiasing mechanism \cite{owens2024multi}; we therefore treat the sign of the C3 effect as an empirical question and report it regardless of direction.

Statistical treatment is paired across counterfactual variants, with confidence intervals on each disparity metric and corrections for multiple comparisons across the configuration × cue-channel grid; non-significant results will be reported as such rather than suppressed.

\subsection{6.4 The Role This Study Plays}

This audit is intentionally modest in scale and ambitious only in framing. Its purpose is threefold. First, it provides a \emph{constructive existence proof}: if even one configuration confirms H1, the measurement gap is not merely a conceptual argument but a measured phenomenon, motivating the action-level metrics of Section 4.3 and the open problems of Section 8.1. Second, it serves as a methodological template — pinned models, synthetic counterfactuals, paired CFR/MASD-plus-action metrics — that can be scaled across domains and architectures. Third, it functions as a teaser for the companion \textbf{acting-agent fairness benchmark} that Section 8.1 calls for: the audit demonstrates, in miniature, the instrumentation that a full benchmark would standardize. We are explicit that this study is \emph{not} the survey's main contribution; its weight is evidentiary, anchoring the taxonomy's central thesis with original data rather than carrying the survey's novelty.

We also note its publication role candidly: because the audit pairs a literature synthesis with original, reproducible experimental results, the work is positioned to clear arXiv moderation as a \emph{survey with original results} rather than as a pure review.

\subsection{6.5 Results (to be completed before camera-ready)}

\emph{The audit has been fully specified but not yet executed. Numerical results — CFR, MASD, and action-level disparity (tool-selection, decision, and escalation rates) for configurations C0–C3 across the demographic cue channels, with confidence intervals and significance tests — will be reported here in the camera-ready version. We will release the synthetic profile set, prompts, pinned model identifiers, and analysis code to enable exact reproduction. No results are claimed at this stage; the placeholder is intentional, and any reading of Hypotheses H1–H2 as confirmed should await the completed run.}



\section{7. Coverage Map & Cross-Analysis}

Having traced where unfairness enters each agent component (§3), measured it (§4), and mitigated it (§5), we now project the corpus back onto its two organizing axes simultaneously. Figure~\ref{fig:coverage} renders the coverage matrix: the six agent components (rows) against the three formal fairness dimensions plus an \emph{unspecified} column (columns), where each cell counts corpus papers that engage that component under that framing. Table~0 fixes the survey's position against its nearest neighbors. Read together, the two visuals make the survey's central empirical claim falsifiable: the field is dense exactly where it is easy and empty exactly where agentic deployment makes fairness consequential.

\subsection{7.1 The counterfactual column is essentially empty}

The headline of Figure~\ref{fig:coverage} is a near-vacant column. Across all six components and the 57 component-tagged corpus papers, \textbf{formal counterfactual fairness is the named framing in exactly one paper} — \texttt{she2025fairsense}, under long-horizon drift, which simulates demographic-conditioned trajectories of an ML-enabled system over time. Tool/API selection, memory \& retrieval, multi-agent delegation, planning/decomposition, and user modeling/personalization register \textbf{zero} counterfactual entries each. This is the survey's sharpest finding, and it is a structural inversion of need. The counterfactual flip — re-running the \emph{same} agent on a profile altered only in a protected attribute and asking whether the trajectory changed — is the natural operationalization of the agentic fairness question, "would the agent have \emph{acted} differently for another demographic?" Group and individual metrics, computed over output distributions, cannot answer it directly; counterfactual flip can. The dimension that agentic harm most requires is the dimension the field has least formalized. We note this is a corpus-relative claim: the \emph{methods} exist — CFR/MASD in \texttt{mayilvaghanan2026counterfactual}, the counterfactual audit of \texttt{amirimargavi2026equal}, intervention-consistency in \texttt{basu2026names}, and the CAFFE/SCOPE counterfactual datasets (\texttt{parziale2025systematic}, \texttt{parziale2026scope}) — but they sit in §4's QA-level benchmark inventory, not yet wired through any agent component. The gap is one of application, not invention.

\subsection{7.2 Group fairness clusters only where decisions are explicit}

Where a formal dimension \emph{is} invoked, it is overwhelmingly group fairness, and it clusters in precisely two components: \textbf{multi-agent delegation (5 group-framed papers, 1 individual)} and \textbf{planning/decomposition (4 group, 1 individual)}. These are the components where the agent emits an explicit, attributable decision — a role assignment, a debate verdict, a task allocation, a discrete action — against which selection-rate or outcome-parity gaps can be tabulated. Multi-agent work (\texttt{madigan2025emergent}, \texttt{nguyen2025social}, \texttt{li2025from}, \texttt{mirza2025malibu}) measures disparity across agent personas and emergent group outcomes; planning work (\texttt{li2025actions}, \texttt{hall2025guiding}, \texttt{parziale2026once}) measures decision and allocation disparity. The lesson is structural: formalization follows decision explicitness. Components that surface a clean decision variable attract group metrics; components whose harm is diffuse do not.

\subsection{7.3 Tool, memory, and personalization harms are documented but rarely formalized}

The complementary pattern is the heavily weighted \emph{unspecified} column. Three components are dominated by it almost entirely: \textbf{memory \& retrieval (0 formal, 12 unspecified)}, \textbf{tool/API selection (1 group, 11 unspecified)}, and \textbf{user modeling/personalization (1 group, 11 unspecified)}. Memory \& retrieval is the most striking — \textbf{not a single paper} in this branch names a formal fairness dimension, despite a substantial body documenting real harm: RAG that degrades fairness even for vigilant users (\texttt{hu2024free}), multi-turn fairness erosion (\texttt{fan2024fairmt}), and belief drift from accumulating context (\texttt{geng2025accumulating}). These papers \emph{demonstrate} disparity convincingly; they do not \emph{measure} it against group, individual, or counterfactual criteria, reporting instead bespoke or qualitative gaps. The same holds for tool selection (\texttt{blankenstein2025biasbusters}, \texttt{sneh2025tooltweak}) and personalization (\texttt{kantharuban2024stereotype}, \texttt{pawar2025presumed}, \texttt{mire2025rejected}). Harm in these components is real and replicated, but it lives outside the formal-fairness vocabulary — which means it cannot be benchmarked, tracked across versions, or compared across systems. Documentation without formalization is the field's modal state for half the agent pipeline.

\subsection{7.4 Synthesis and differentiation}

Three regularities fall out of the matrix. First, \textbf{formalization tracks decision explicitness, not harm severity} — components emit formal metrics in proportion to how cleanly they expose a decision variable, regardless of where the deployment stakes actually lie. Second, \textbf{the formal vocabulary thins toward the counterfactual} exactly as the question shifts from "what did the agent output?" to "would it have acted differently?" — the agentic question. Third, \textbf{the modal corpus paper documents disparity without naming a dimension} (51 of the 57 component-tagged entries sit in the \emph{unspecified} column), leaving the field rich in evidence and poor in commensurable measurement.

This map is also what distinguishes the survey from its neighbors (Table~0). Where \texttt{chu2024fairness} organizes single-turn LLM fairness by metric and \texttt{gallegos2024bias} by datasets/mitigation, neither resolves \emph{where in an acting pipeline} a harm originates. Where \texttt{mohammadi2025evaluation} treats fairness as one sub-topic of agent evaluation and \texttt{vatsal2026agentic} as one of seven clinical dimensions, neither makes component coverage its unit of analysis. Where \texttt{mayilvaghanan2026counterfactual} supplies the counterfactual method we generalize, it does so for a single domain. By plotting the corpus on the component × dimension grid, this survey turns the measurement gap from an assertion into a coordinate — the empty counterfactual column, the unspecified-dominated rows — and hands §8 a concrete target: an \emph{acting-agent} benchmark must populate the cells this matrix leaves blank.



\section{8. Open Problems & Research Agenda}

The preceding taxonomy exposes a field whose investigative apparatus is systematically misaligned with its object of study. Bias is documented at every agent component — tool selection \cite{blankenstein2025biasbusters}, retrieval \cite{wu2024rag}, delegation \cite{nguyen2025social}, planning \cite{li2025actions}, personalization \cite{kantharuban2024stereotype}, and long-horizon drift \cite{ma2026implicit} — yet the measurement instruments that certify a system as "fair" almost uniformly probe single-turn question answering. This section converts that mismatch into a concrete agenda. We frame six open problems; for each we state why it matters for consequential deployment and what a credible solution must deliver. The through-line is that fairness of \emph{acting} agents is under-measured, under-governed, and under-mitigated relative to the autonomy these systems already exercise in hiring, lending, healthcare, and public benefits.

\subsection{8.1 The measurement gap: benchmarks for acting agents}

\textbf{The problem.} The dominant fairness benchmarks score answers, not actions. Clinical fairness sets exemplify the ceiling of current practice and its limit: FairMedQA evaluates demographic bias in medical question answering \cite{xiao2025fairmedqa}; mFARM assembles multi-faceted, harm-grounded fairness assessment for clinical decision support but operates on static decision vignettes rather than executed care trajectories \cite{adappanavar2025mfarm}; MedEqualQA isolates bias through counterfactual demographic edits to otherwise identical questions \cite{ghosh2025medequalqa}; and equity toolboxes for health LLMs surface harms at the level of generated text \cite{pfohl2024toolbox}. Each is methodologically careful, and each stops at the QA boundary. None measures what a clinical \emph{agent} does when it selects which guideline tool to invoke, which prior note to retrieve, whether to escalate to a human, or how it sequences a multi-step workup — precisely the loci where the taxonomy of Section 3 locates the agentic harms. The few works that score actions rather than utterances confirm the divergence: agent decisions reveal implicit biases that token-level probes miss \cite{li2025actions}, and decomposing unfairness in a tool-using chest X-ray agent attributes disparity to specific pipeline stages invisible to answer-level scoring \cite{xu2026ducx}.

\textbf{Why it matters.} A system can be near-parity on FairMedQA and still triage, escalate, or order differently across demographic counterfactuals once it acts, because disparity enters through tool routing and retrieval rather than the final string. Certifying agents with QA benchmarks therefore licenses deployment on evidence that does not bind the deployed behavior.

\textbf{What a solution needs.} An acting-agent fairness benchmark must (i) hold counterfactual demographic profiles fixed across full trajectories rather than single prompts, generalizing the counterfactual flip-rate (CFR) and mean-absolute-score-difference (MASD) machinery of \cite{mayilvaghanan2026counterfactual} from a single customer-service turn to multi-step execution; (ii) instrument \emph{action-level} disparity — differential tool invocation, retrieval composition, decision, and escalation rate — not just output text; and (iii) ground tasks in a domain where actions carry real stakes. We are developing \textbf{FairMedAgent}, a forthcoming companion benchmark of this kind for clinical acting agents, which evaluates tool-using and multi-step clinical agents under matched demographic counterfactuals and reports action-level disparities alongside outcome parity. FairMedAgent is the executable counterpart to the QA-level clinical sets above and the empirical instrument the gap demands; the cross-framework audit of Section 6 is its proof-of-concept.

\subsection{8.2 Long-horizon and compounding fairness}

\textbf{The problem.} Fairness is overwhelmingly evaluated as a one-shot property, but agents persist — accumulating memory, conditioning on their own prior actions, and operating inside feedback loops. Bias compounds along these trajectories. Implicit bias in long-term memory has been shown to accumulate and propagate, so that a small initial skew is reinforced rather than averaged out across interactions \cite{ma2026implicit}. The sequential-decision literature establishes the formal stakes: static, per-step fairness constraints fail to deliver equitable long-term outcomes, and achieving long-term fairness requires reasoning about how present decisions shape future eligibility and state \cite{hu2022achieving}. Adapting static fairness to sequential decision-making shows that naive transfer of one-shot metrics can degrade equal long-term benefit rates, motivating mitigation that targets cumulative rather than instantaneous disparity \cite{xu2024adapting}.

\textbf{Why it matters.} Agentic deployments in benefits adjudication, longitudinal care, and credit are inherently sequential; a system that is fair at every isolated step can still drive divergent cumulative outcomes across groups. One-shot certification is blind to exactly the dynamics that determine who is ultimately served.

\textbf{What a solution needs.} Longitudinal evaluation protocols that measure disparity over trajectories, not turns — tracking how memory state, self-conditioning, and feedback modify group outcomes over time — together with long-horizon mitigation objectives (e.g., equal long-term benefit rate \cite{xu2024adapting}) and explicit models of the induced state dynamics \cite{hu2022achieving}.

\subsection{8.3 Multi-agent fairness dynamics}

\textbf{The problem.} Multi-agent orchestration introduces failure modes with no single-agent analog: bias can \emph{emerge} from interaction even among individually well-behaved agents. Emergent bias and unfairness arise in multi-agent decision systems as a property of the collective rather than any constituent model \cite{madigan2025emergent}, and stereotypical bias has been shown to emerge, propagate, and \emph{amplify} through agent-to-agent communication, so that the social cost of a multi-agent system exceeds the sum of its parts \cite{nguyen2025social}. Disparity here is a function of topology, role assignment, and message-passing dynamics, none of which a per-agent fairness audit can observe.

\textbf{Why it matters.} Production agent systems are increasingly multi-agent (debate, critic, delegation, role specialization). If amplification is emergent, then component-wise certification of each agent provides no guarantee about the system, and the very architectures marketed as more capable may be more unfair.

\textbf{What a solution needs.} Fairness metrics defined over the multi-agent \emph{system} — quantifying emergence and amplification as a function of communication topology and role structure \cite{madigan2025emergent,nguyen2025social} — plus interventions that act on the interaction graph rather than on isolated agents (Section 5.4).

\subsection{8.4 Intersectional and sociotechnical harms in deployment}

\textbf{The problem.} Bias research frequently isolates one protected attribute on a clean benchmark, but real deployments harm people at intersections and through sociotechnical channels that single-axis, lab-clean evaluation cannot surface. Large-scale evidence from deployed-adjacent settings is the most damning: sociodemographic biases in medical decision making by LLMs produce differential management recommendations across patient demographics at clinical scale \cite{omar2025sociodemographic}; a mortgage-underwriting experiment finds measurable racial disparities in LLM lending decisions \cite{bowen2024measuring}; and a judicial-fairness evaluation documents systematic unfairness in legal decision contexts \cite{hu2025llms}. These harms are jointly demographic and contextual — they depend on the deployment's stakes, its decision structure, and the combination of attributes a real applicant carries.

\textbf{Why it matters.} High-stakes domains are exactly where intersectional and contextual harms concentrate and where the cost of missing them is borne by the most vulnerable. A survey that organized fairness by metric would file these as scattered case studies; organizing by agent component shows they are the downstream signature of component-level disparities measured under realistic, multi-attribute conditions.

\textbf{What a solution needs.} Evaluation that is intersectional by construction (joint, not marginal, attribute combinations), domain-grounded in the decision structure of the target setting \cite{omar2025sociodemographic,bowen2024measuring,hu2025llms}, and attentive to the sociotechnical context — who is affected, what recourse exists, and how the agent's autonomy changes the harm surface.

\subsection{8.5 Standardization, auditing, and governance}

\textbf{The problem.} There is no shared standard for what it means to audit an LLM agent for fairness, nor an agreed reporting format for action-level disparity. The field has converging \emph{ingredients} — counterfactual frameworks such as CAFFE \cite{parziale2025systematic}, stereotyped-prompt datasets such as SCOPE \cite{parziale2026scope}, use-case-specific evaluation tooling \cite{bouchard2024bring}, intervention-consistency tests \cite{basu2026names}, and dataset surveys cataloguing the landscape \cite{zhang2025datasets} — but no standardization that turns them into a repeatable, comparable audit of an \emph{acting} system. Broader agent-evaluation and trustworthy-agent surveys treat fairness as one dimension among many \cite{mohammadi2025evaluation,yu2025survey,vatsal2026agentic}, leaving fairness auditing of agents without a dedicated protocol.

\textbf{Why it matters.} Agents are being deployed into hiring, lending, healthcare, and public-benefits pipelines — domains of clear national importance — faster than auditable standards exist to govern them. Standardized, action-level fairness auditing is a precondition for accountable deployment and for any regulatory or procurement regime that would condition use on demonstrated equity; the absence of such standards is itself a public-interest gap this line of work is positioned to address.

\textbf{What a solution needs.} A reference auditing protocol specifying trajectory-level counterfactual construction, action-level disparity metrics, and standardized reporting; reusable, domain-instantiable instruments (extending \cite{parziale2025systematic,parziale2026scope,bouchard2024bring}); and governance hooks — provenance, logging, and recourse — that make agent decisions contestable after the fact.

\subsection{8.6 Mitigation at the agent level, not the token level}

\textbf{The problem.} The mitigation literature overwhelmingly intervenes on the model's \emph{outputs} — debiasing prompts \cite{furniturewala2024thinking,li2024prompting}, preference optimization and alignment \cite{allam2024biasdpo,bai2022constitutional}, inference-time activation steering \cite{li2025fairsteer}, and output guardrails \cite{dong2024building}. These operate on tokens. But Section 3 locates agentic harm in \emph{actions and architecture}: which tool is chosen, what is retrieved into context, how roles are assigned, how a plan is decomposed. Token-level debiasing leaves these channels untouched. The few process-aware interventions point the way — bias-mitigation agents that optimize source selection for balanced retrieval \cite{singh2025biasb}, bias-aware retrieval agents \cite{singh2025bias}, identity anonymization to neutralize role-induced debate bias \cite{choi2025identity}, and multi-agent debiasing frameworks \cite{owens2024multi,xu2024mitigating} — but they remain scattered relative to the dominant token-level paradigm.

\textbf{Why it matters.} If disparity enters through tool routing and retrieval composition, an output-fair model can still act unfairly; conversely, constraining the agent's \emph{process} can reduce disparity even where the base model is biased. Mitigation must meet harm at the layer where it originates.

\textbf{What a solution needs.} Architectural and process-level interventions keyed to specific components — tool-selection constraints, retrieval-composition controls, role and identity hygiene in delegation, and planning audits — evaluated by the same action-level disparity metrics called for in 8.1, so that mitigation is credited for changing what the agent \emph{does}, not merely what it says.


\textbf{Synthesis.} These six problems share one root: the autonomy that makes agents valuable — acting, persisting, delegating, and self-conditioning — is precisely what current QA-level measurement, single-shot certification, and token-level mitigation cannot see. Closing the gap requires instruments, protocols, and interventions defined over \emph{trajectories and actions}. The remainder of this agenda — and the forthcoming FairMedAgent benchmark — is an attempt to build them.



\section{Conclusion}

The case for analyzing fairness \textbf{by agent component} rests on a structural observation: LLM-based agents do not merely answer questions — they select tools, accumulate memories across turns, delegate tasks across role-assigned subagents, decompose goals into multi-step plans, adapt to inferred user identities, and persist behavior over long horizons \cite{wang2024survey,xi2023rise}. Each of these pipeline stages constitutes an independent locus at which bias can enter, amplify, and compound into consequential harm \cite{gallegos2024bias,madigan2025emergent}. A QA-level fairness score, averaged over outputs, is blind to this architecture: it can report a neutral aggregate while tool-selection disparity \cite{blankenstein2025biasbusters,xu2026ducx}, retrieval poisoning \cite{wu2024rag,hu2024free}, persona-induced role bias \cite{cao2026from,gupta2023bias}, planning-stage stereotyping \cite{li2025actions,parziale2026once}, personalization harms \cite{weissburg2024llms,kantharuban2024stereotype}, and multi-turn drift \cite{ma2026implicit,simhi2026old} each operate undetected beneath the surface.

This survey has presented the first dedicated taxonomy of fairness in LLM-based agents organized along this component axis — spanning tool and API selection (§3.1), memory and retrieval (§3.2), multi-agent delegation and role assignment (§3.3), planning and decomposition (§3.4), user modeling and personalization (§3.5), and long-horizon trajectory drift (§3.6), with a cross-cutting account of how component-level harms compound into an agentic multiplier (§3.7). Against this architecture we mapped the evaluation landscape (§4), characterizing the family of group, individual, and counterfactual metrics and generalizing the CFR and MASD methods of \citet{mayilvaghanan2026counterfactual} across all component types, while Table~2 documents that nearly every existing benchmark — clinical fairness sets \cite{xiao2025fairmedqa,adappanavar2025mfarm,ghosh2025medequalqa}, hiring and lending audits \cite{an2025measuring,azime2025accept}, compositional evaluation suites \cite{wang2024ceb} — scores model outputs rather than agent actions. The mitigation landscape (§5) follows the same component logic, distinguishing prompting-level interventions \cite{furniturewala2024thinking,li2024prompting} from alignment-stage approaches \cite{allam2024biasdpo,bai2022constitutional}, architectural guardrails \cite{dong2024building}, and the emerging class of process-level and multi-agent debiasing strategies \cite{owens2024multi,choi2025identity,ki2025multiple}. The original cross-framework audit (§6) confirmed empirically what the coverage matrix (§7) shows structurally: action-level disparity — which tools an agent invokes, which decisions it escalates, which candidates it advances — is real, measurable, and invisible to today's QA-level scoring.

The measurement gap is therefore not merely a gap in tooling; it is a gap in accountability. Agents are already deployed in hiring \cite{veldanda2023emily,wilson2024gender,nghiem2024you}, clinical decision support \cite{poulain2024bias,omar2025sociodemographic,benkirane2024diagnose}, lending \cite{bowen2024measuring,feng2023empowering}, and education \cite{weissburg2024llms,yang2025prompt} — consequential domains where differential treatment across race, gender, socioeconomic status, dialect, and their intersections carries direct legal and ethical weight \cite{bommasani2021opportunities,bender2021dangers}. The national-importance stakes are not speculative: fairness failures that compound silently across agent trajectories \cite{she2025fairsense,hu2022achieving} can systematically close access to employment, credit, and care for protected groups at a scale and opacity no manual audit is equipped to catch.

The immediate priority the agenda of §8 identifies is therefore benchmark construction for \emph{acting} agents — evaluation infrastructure that scores trajectories and decisions rather than tokens and outputs, that applies counterfactual demographic perturbations at the action level \cite{parziale2025systematic,amirimargavi2026equal,basu2026names}, and that covers the component-level granularity the taxonomy exposes. Closing this gap is the precondition for every subsequent advance in mitigation, governance, and auditing of agentic AI. The field has the conceptual vocabulary; it now needs the instruments.

\bibliographystyle{ACM-Reference-Format}
\bibliography{references}

\end{document}
